
\documentclass{article}
\usepackage{colortbl}
\usepackage{makecell}
\usepackage{multirow}
\usepackage{supertabular}

\begin{document}

\newcounter{utterance}

\centering \large Interaction Transcript for game `cladder', experiment `full\_v1.5\_default', episode 7517 with refineonly{-}Qwen3.5{-}4B.
\vspace{24pt}

{ \footnotesize  \setcounter{utterance}{1}
\setlength{\tabcolsep}{0pt}
\begin{supertabular}{c@{$\;$}|p{.15\linewidth}@{}p{.15\linewidth}p{.15\linewidth}p{.15\linewidth}p{.15\linewidth}p{.15\linewidth}}
    \# & \multicolumn{2}{c}{Player} && \multicolumn{2}{c}{Game Master} \\
    \hline

    \theutterance \stepcounter{utterance}  
    & & & \multicolumn{4}{p{0.6\linewidth}}{
        \cellcolor[rgb]{0.9,0.9,0.9}{
            \makecell[{{p{\linewidth}}}]{
                \texttt{\tiny{[P1$\langle$GM]}}
                \texttt{You are an expert in causal inference. The following question is not a typical commonsense query, but rather a meticulously designed question created by a professor specializing in causal inference, intended to assess the students' mastery of the course content.} \\
\texttt{Imagine a self{-}contained, hypothetical world with only the following conditions, and without any unmentioned factors or causal relationships: Wibl has a direct effect on vubr and zuph. Zuph has a direct effect on uvzi. Vubr has a direct effect on uvzi. Method 1: We look at how zuph correlates with uvzi case by case according to vubr. Method 2: We look directly at how zuph correlates with uvzi in general. To understand how zuph affects uvzi, is it more correct to use the Method 1 than Method 2?} \\
\texttt{Start your answer with 'Yes' or 'No', followed by additional reasoning or evidence to support your explanation.} \\
            }
        }
    }
     \\ \\

    \theutterance \stepcounter{utterance}  
    & \multicolumn{4}{p{0.6\linewidth}}{
        \cellcolor[rgb]{0.561,0.737,0.561}{
            \makecell[{{p{\linewidth}}}]{
                \texttt{\tiny{[P1$\rangle$GM]}}
                \texttt{Yes} \\
\\ 
\texttt{**Reasoning:**} \\
\\ 
\texttt{To determine the correct method for understanding the causal effect of \$zuph\$ on \$uvzi\$, we must analyze the causal structure and the presence of confounding variables.} \\
\\ 
\texttt{1.  **Causal Structure Analysis**:} \\
\texttt{    *   \$wibl \textbackslash{}rightarrow zuph\$} \\
\texttt{    *   \$wibl \textbackslash{}rightarrow vubr\$} \\
\texttt{    *   \$vubr \textbackslash{}rightarrow uvzi\$} \\
\texttt{    *   \$zuph \textbackslash{}rightarrow uvzi\$} \\
\\ 
\texttt{    In this diagram, \$wibl\$ is a **common cause** (a confounder) of both \$zuph\$ and \$vubr\$. Furthermore, \$vubr\$ acts as a **mediator** on the path from \$wibl\$ to \$uvzi\$. The critical issue here is that \$vubr\$ is a descendant of the confounder \$wibl\$.} \\
\\ 
\texttt{2.  **Evaluating Method 2 (General Correlation)**:} \\
\texttt{    Method 2 looks at the raw correlation between \$zuph\$ and \$uvzi\$. In this unadjusted setting, the observed association between \$zuph\$ and \$uvzi\$ is a mixture of:} \\
\texttt{    *   The direct causal effect (\$zuph \textbackslash{}rightarrow uvzi\$).} \\
\texttt{    *   The indirect effect mediated through \$vubr\$ (\$zuph \textbackslash{}leftarrow wibl \textbackslash{}rightarrow vubr \textbackslash{}rightarrow uvzi\$).} \\
\\ 
\texttt{    Because \$wibl\$ causes both \$zuph\$ and \$vubr\$, and \$vubr\$ causes \$uvzi\$, there is a spurious correlation induced by \$wibl\$. Specifically, high values of \$wibl\$ lead to high \$zuph\$ and high \$vubr\$, which in turn leads to high \$uvzi\$. This creates a positive association between \$zuph\$ and \$uvzi\$ that is partially (or wholly, depending on the strength of the direct vs. indirect paths) driven by the common cause \$wibl\$, not the direct causal link. Therefore, Method 2 yields a biased estimate of the direct causal effect.} \\
\\ 
\texttt{3.  **Evaluating Method 1 (Stratification/Conditioning)**:} \\
\texttt{    Method 1 looks at the correlation between \$zuph\$ and \$uvzi\$ "case by case according to \$vubr\$." This implies conditioning on \$vubr\$ (stratifying the data by levels of \$vubr\$).} \\
\\ 
\texttt{    *   If we condition on \$vubr\$, we are effectively comparing cases where \$vubr\$ is held constant.} \\
\texttt{    *   By fixing \$vubr\$, we block the path \$wibl \textbackslash{}rightarrow vubr \textbackslash{}rightarrow uvzi\$. Since \$vubr\$ is a descendant of the confounder \$wibl\$, conditioning on it blocks the backdoor path from \$zuph\$ to \$uvzi\$ that flows through \$} \\
            }
        }
    }
    & & \\ \\

    \theutterance \stepcounter{utterance}  
    & & & \multicolumn{2}{p{0.3\linewidth}}{
        \cellcolor[rgb]{0.9,0.9,0.9}{
            \makecell[{{p{\linewidth}}}]{
                \texttt{\tiny{[GM$|$GM]}}
                \texttt{Yes} \\
            }
        }
    }
    & & \\ \\

    \theutterance \stepcounter{utterance}  
    & & & \multicolumn{2}{p{0.3\linewidth}}{
        \cellcolor[rgb]{0.9,0.9,0.9}{
            \makecell[{{p{\linewidth}}}]{
                \texttt{\tiny{[GM$|$GM]}}
                \texttt{yes} \\
            }
        }
    }
    & & \\ \\

    \theutterance \stepcounter{utterance}  
    & & & \multicolumn{2}{p{0.3\linewidth}}{
        \cellcolor[rgb]{0.9,0.9,0.9}{
            \makecell[{{p{\linewidth}}}]{
                \texttt{\tiny{[GM$|$GM]}}
                \texttt{game\_result = WIN} \\
            }
        }
    }
    & & \\ \\

\end{supertabular}
}

\end{document}
